\documentclass[12pt]{article}
\usepackage[utf8]{inputenc}
\usepackage[T1]{fontenc}
\usepackage{hyperref}
\title{Baiyamo Language Resources for K-12 Curriculum Development}
\author{Compiled by Claude Sonnet 5.5 (claude-sonnet-5-5) for Andr\'as Kornai}
\date{}
\begin{document}
\maketitle

\section{Introduction}
Baiyamo (Papi, Paupe; ISO 639-3 ppe) is a small language of Papua New Guinea spoken in one village on the lower Frieda River south of the Upper Sepik, with a few speakers in a hamlet on the Leonhard Schultze River. Its classification is uncertain (Sepik by some, unclassified in Glottolog). Estimates of speakers vary, about 70 in the reference description and 220 in Ethnologue 2011. All speakers also know Tok Pisin.

\section{Primary Educational Resources}
No primers, textbooks or school programmes were found.

\section{Linguistic and Reference Materials}
\begin{itemize}
\item Overview article: \url{https://en.wikipedia.org/wiki/Papi_language}
\item Numeral data page (numerals 1--4 only): \url{https://lingweb.eva.mpg.de/channumerals/Baiyamo.htm}
\end{itemize}
No dictionary or grammar was found.

\section{Religious and Community Resources}
None found.

\section{Digital and Audio Resources}
\begin{itemize}
\item Endangered Languages Project entry: \url{https://endangeredlanguages.com/lang/2645}
\end{itemize}

\section{Curriculum Development Context}
Documentation is minimal, so a curriculum would require new fieldwork and materials developed with the community.
\end{document}